\documentclass[12pt]{article}

\usepackage[portuguese]{babel}
\usepackage{float}
\usepackage{listings}
\usepackage{indentfirst}
\usepackage{amsfonts, amssymb, amsmath}
\usepackage{graphicx}
\usepackage{tikz}
\usepackage{enumitem}
\usepackage[colorlinks=true, allcolors=blue]{hyperref}
\usepackage[letterpaper,
            top=1.5cm,
            bottom=1.5cm,
            left=2.75cm,
            right=2.75cm,
            marginparwidth=1.75cm]{geometry}

\title{Psicologia no Trabalho}
\author{Prof$\mathrm{^a}$ Tamara Menezes}
\date{\today}

\begin{document}
\maketitle
\tableofcontents

\section{Trabalho e Subjetividade}
\begin{itemize}
    \item Trabalho como atividade fundamental na construção do sujeito 
    $\rightarrow$ investimento afetivo
    \item Espaço de edificação de sentido
    \begin{itemize}
        \item Identidade pessoal e social
    \end{itemize}
\end{itemize}
\subsection{Subjetividade}
\begin{itemize}
    \item Experiencias individuais, significados que permeiam a vida de uma 
    pessoa
    \begin{itemize}
        \item Exemplos: história de vida, linguagem, cultura, familia \ldots
    \end{itemize}
    \item Modos de ser e de expressão de um sujeito no mundo
    \item Expressa em lugares de saner e produz a verdade de um sujeito
    
    Genero, raça e classe social não são indissociaveis
\end{itemize}

\section{A Dupla Face da Vida Organizacional}
\begin{itemize}
    \item O trabalho é um fenomeno social e humano. Para entender a organização,
    precisamos de duas lentes:
    \begin{enumerate}
        \item Como o individuo sente e age
        \item Como o coletivo contrói significados
    \end{enumerate}
    \begin{tikzpicture}
        \draw (0,0) circle (3cm) node[below left] {};
        \draw (3.5,0) circle (3cm) node[below right] {};

        \node at (-0.5,0) {10};
        \node at (4,0) {15};
        \node at (1.5,0) {5};
        \end{tikzpicture}
\end{itemize}
    \subsection{O trabalhar}
    \begin{itemize}
        \item Trabalho prescrito (Manuais, regras, previsões)
        \item Trabalho real (Imprevistos, falhas, dinâmicas)
        \item Trabalhar é o que o sujeito acrescenta de si mesmo para fazer o 
        sistema funcionar quando a regra falha. O trabalho efetivo é, em sua 
        essência, invisível
    \end{itemize}

\end{document}